%! Planning a Study — Unit 3, Lesson 3.1 — Slides (Beamer)
\documentclass[aspectratio=169, 11pt]{beamer}
\usepackage{apstats-beamer}   % provides \navyheader{} and \sectionlabel[color]{}

\begin{document}

% TITLE SLIDE (CourseName is not defined in beamer, so the name is literal)
{
\setbeamercolor{background canvas}{bg=navy}
\begin{frame}[plain]
  \vspace{2.8cm}\hspace{0.55cm}%
  \begin{minipage}{0.9\textwidth}
    {\color{goldacc}\bfseries\small LESSON 3.1}\\[0.5cm]
    {\color{white}\bfseries\LARGE Planning a Study}\\[1.2cm]
    {\color{skymid}\small Statistics~$\cdot$~Unit 3}
  \end{minipage}
\end{frame}
}

% GRADUAL-RELEASE deck: targets → warm-up → I DO divider → one frame per instruction row of the
% notes (Ask How; Sample; Study Type, which carries the crux) → WE DO (The Study App) → spoken
% debrief → close & START THE HOMEWORK, which is the individual practice.
% There is no group activity and no independent practice set in the notes.
\newcommand{\redink}[1]{{\color{redacc}\bfseries #1}}

% ── LEARNING TARGETS ─────────────────────────────────────────────────────────
\begin{frame}[t, plain]
  \navyheader{Today's targets}
  \sectionlabel{Lesson 3.1}
  By the end of today, I can\ldots
  \begin{itemize}\setlength{\itemsep}{4pt}
    \item tell a \textbf{chance-based method} from a \textbf{non-chance method}.
    \item name the \textbf{population} and the \textbf{sample}, say what a \textbf{census} is,
          and generalize only to the population sampled.
    \item classify a study as an \textbf{observational study} or an \textbf{experiment}.
    \item explain why only a well-designed experiment can show \textbf{cause and effect}.
  \end{itemize}
  \begin{block}{How today works}
    Warm-up $\to$ \textbf{notes together} $\to$ \textbf{one we work as a class} $\to$
    debrief $\to$ \textbf{you start the homework here, alone}.
  \end{block}
\end{frame}

% ── WARM-UP ──────────────────────────────────────────────────────────────────
\begin{frame}[t, plain]
  \navyheader{Warm-up}
  \sectionlabel{5 minutes $\cdot$ on your own}
  \begin{enumerate}\setlength{\itemsep}{8pt}
    \item A news site's online poll: 2{,}400 students respond, \textbf{85\%} want a later start.
          Who answers a poll like this? Does 85\% describe \emph{all} students?
    \item 900 students, one lunch question. One reason to ask all 900 --- and one reason to ask
          only some.
    \item More sleep, higher test scores. Explanatory? Response? Does it show \emph{cause}?
  \end{enumerate}
  \vspace{6pt}
  \begin{block}{Each one comes back}
    Item 1 is Row 1 of the notes, item 2 is Row 2, item 3 is Row 3.
  \end{block}
\end{frame}

% ── I DO — direct instruction begins ─────────────────────────────────────────
{
\setbeamercolor{background canvas}{bg=navy}
\begin{frame}[plain]
  \vspace{1.8cm}\hspace{0.8cm}%
  \begin{minipage}{0.86\textwidth}
    {\color{goldacc}\bfseries\small GUIDED NOTES \ \textbullet\ \ I DO}\\[0.7cm]
    {\color{white}\bfseries\Large Open to your Guided Notes. Fill each blank as we name it
    --- not before.}\\[0.9cm]
    {\color{skymid}\small Four terms go in the vocabulary box today: population, sample,
    observational study, experiment.}
  \end{minipage}
\end{frame}
}

% ── ROW 1: ASK HOW ───────────────────────────────────────────────────────────
\begin{frame}[t, plain]
  \navyheader{Before you trust a number, ask how}
  \sectionlabel{notes, Ask How $\cdot$ problems 1--3}
  \vspace{2pt}
  \begin{columns}[t]
    \begin{column}{0.48\textwidth}
      \begin{block}{Chance-based method}
        A \textbf{random} process decides who is in the data.
      \end{block}
      \vspace{3pt}
      \begin{block}{Non-chance method}
        \textbf{People} decide --- the researcher takes whoever is handy, or respondents
        volunteer.
      \end{block}
    \end{column}
    \begin{column}{0.48\textwidth}
      Methods that do not rely on chance give \redink{untrustworthy} conclusions.
      \vspace{8pt}

      {\small\textbf{Who decided --- chance or people?}}
      \begin{enumerate}\setlength{\itemsep}{3pt}\small
        \item A radio call-in vote on a new stadium.
        \item A coach asks his varsity players about sports spending.
        \item A computer picks 60 IDs from the full roster.
      \end{enumerate}
    \end{column}
  \end{columns}
\end{frame}

% ── ROW 2: POPULATION, SAMPLE, CENSUS ────────────────────────────────────────
\begin{frame}[t, plain]
  \navyheader{Who we want to know about, and who we ask}
  \sectionlabel{notes, Sample $\cdot$ problems 4--6}
  \vspace{2pt}
  \begin{columns}[t]
    \begin{column}{0.48\textwidth}
      \begin{block}{Population}
        \textbf{All} the individuals we want to know about.
      \end{block}
      \vspace{3pt}
      \begin{block}{Sample}
        The \textbf{subset} of the population we actually measure.
      \end{block}
      \vspace{3pt}
      \begin{block}{Census}
        Measures \textbf{every} individual in the population.
      \end{block}
    \end{column}
    \begin{column}{0.48\textwidth}
      A principal asks about weekly homework hours. A computer picks \textbf{100} of her
      school's \textbf{1{,}200} students at random.
      \vspace{8pt}

      Population? Sample? Why not a census?
      \vspace{8pt}

      A sample speaks only for \redink{the population it came from} --- her school, not the
      state.
    \end{column}
  \end{columns}
\end{frame}

% ── ROW 3: STUDY TYPE — THE CRUX ─────────────────────────────────────────────
\begin{frame}[t, plain]
  \navyheader{Watching or doing?}
  \sectionlabel{notes, Study Type $\cdot$ problems 7--10 $\cdot$ the whole point of today}
  \vspace{2pt}
  \begin{columns}[t]
    \begin{column}{0.46\textwidth}
      \begin{block}{Observational study}
        Observe and measure. \textbf{No treatment imposed.} A sample survey is one kind.
        Shows only an \textbf{association}.
      \end{block}
    \end{column}
    \begin{column}{0.46\textwidth}
      \begin{block}{Experiment}
        Researchers \textbf{assign treatments} and measure the response. Only a well-designed
        experiment can show \textbf{cause}.
      \end{block}
    \end{column}
  \end{columns}
  \vspace{8pt}
  A sleep lab fits 50 adults with wrist monitors and records their exact sleep and caffeine for
  a month. ``It was measured so carefully that it is an experiment.'' \emph{Agree?}
  \vspace{6pt}
  \begin{block}{The one question that decides it}
    Did the researchers \redink{impose a treatment}? Careful measuring is still watching.
  \end{block}
\end{frame}

% ── WE DO ────────────────────────────────────────────────────────────────────
\begin{frame}[t, plain]
  \navyheader{Guided Practice: The Study App}
  \sectionlabel[goldacc]{We do $\cdot$ you hold the pen $\cdot$ problems 11--14}
  A university wants to know about the study habits of \emph{all} \textbf{20{,}000} of its
  students.
  \vspace{4pt}
  \begin{itemize}\setlength{\itemsep}{3pt}
    \item \textbf{Study A:} a computer picks 250 students at random from the registrar's list;
          each reports weekly study hours.
    \item \textbf{Study B:} 120 intro-biology students are randomly assigned to use a flashcard
          app or not; exam scores are compared.
  \end{itemize}
  \vspace{4pt}
  \begin{block}{Ask it every time}
    Who is the population? Did chance choose the sample? Was a treatment
    \redink{imposed}?
  \end{block}
\end{frame}

% ── DEBRIEF ──────────────────────────────────────────────────────────────────
\begin{frame}[t, plain]
  \navyheader{Debrief: three questions to ask any study}
  \sectionlabel{8 minutes $\cdot$ pens down, papers up --- we say these aloud}
  \vspace{6pt}
  \begin{enumerate}\setlength{\itemsep}{6pt}
    \item \textbf{Who decided who is in the data?} Chance, or people?
    \item \textbf{Who does the sample speak for?} Only the population it was drawn from.
    \item \textbf{Was a treatment imposed?} If not, it is an observational study --- an
          association, never proof of cause.
  \end{enumerate}
  \vspace{8pt}
  \begin{block}{Say it out loud}
    Why is ``they measured it so carefully'' the \emph{wrong} reason to call the sleep lab an
    experiment? \quad Measuring carefully is still \redink{watching}.
  \end{block}
\end{frame}

% ── CLOSE & START THE HOMEWORK ──────────────────────────────────────────────
\begin{frame}[t, plain]
  \navyheader{Homework --- start it now}
  \sectionlabel[goldacc]{5 items $\cdot$ scored 1 to 5 $\cdot$ due the first class after two study halls}
  \begin{itemize}\setlength{\itemsep}{5pt}
    \item Part A: Riverside's dog-park survey --- a social-media poll against 400 households
          picked at random. Chance or not? Population and sample? Why not a census?
    \item Part B: one multiple-choice item, with a sentence of justification, and a step-counter
          study the paper calls an experiment.
  \end{itemize}
  \begin{block}{Silent and alone --- I am circulating}
    \emph{Part B, item 2} is the one I am reading over your shoulder. The AP Practice page is
    \textbf{extra credit, scored 1 to 5} --- work it for the reps.
  \end{block}
\end{frame}

% ── CLOSE ────────────────────────────────────────────────────────────────────
{
\setbeamercolor{background canvas}{bg=navy}
\begin{frame}[plain]
  \vspace{1.5cm}\hspace{0.8cm}%
  \begin{minipage}{0.86\textwidth}
    {\color{goldacc}\bfseries\small BEFORE YOU GO}\\[0.7cm]
    {\color{white}\bfseries\Large Ask who decided who is in the data, who the sample speaks for,
    and whether anyone imposed a treatment.}\\[0.9cm]
    {\color{skymid}\small Next: Lesson 3.2 --- Random Sampling Methods: simple random,
    stratified, and cluster samples.}
  \end{minipage}
\end{frame}
}

\end{document}
